\begin{textsample}{POS Dim 3 – ce – Score 16.00 – ce/ce\_inf\_4.txt}  \label{ex:f3_pos_079}
Direitos de aprendizagem e desenvolvimento ( BRASIL, 2017, p. 36 ) Como os direitos estão sendo concretizados na Escola de Educação \textbf{Infantil}?
Conviver com outras \textbf{crianças} e \textbf{adultos}, em \textbf{pequenos} e grandes grupos, utilizando diferentes linguagens, ampliando o conhecimento de si e do outro, o respeito em relação à cultura e às diferenças entre as pessoas. - Diferentes formas de organizar o grupo de \textbf{crianças} têm sido possibilitadas?
Com que frequência?
Há algumas predominantes?
Podem ser incrementadas visando à convivência nos grupos? - Como as \textbf{crianças} interagem nessas diferentes possibilidades?
As \textbf{crianças} expressam suas preferências nesses momentos?
Além da oralidade, outras formas de linguagem das \textbf{crianças} têm sido escutadas e consideradas na organização das situações educativas? - As manifestações culturais que constituem a comunidade escolar têm sido objeto de investigação, elaboração e produção criativa pelas \textbf{crianças}? - A diversidade cultural brasileira e cearense permeia o currículo da Educação \textbf{Infantil}, indo além das datas comemorativas?
Como?
Podem ser ampliadas de forma a enriquecer a construção da identidade comunitária? - As \textbf{crianças} são convidadas a ampliar suas relações de convivência com \textbf{crianças} e \textbf{adultos} que compõem o cotidiano da unidade educativa? \textbf{Brincar} cotidianamente de diversas formas, em diferentes espaços e tempos, com diferentes parceiros ( \textbf{crianças} e \textbf{adultos} ), ampliando e diversificando seu acesso a produções culturais, seus conhecimentos, sua imaginação, sua criatividade, suas experiências emocionais, corporais, \textbf{sensoriais}, expressivas, cognitivas, sociais e relacionais. - O tempo destinado à \textbf{brincadeira} vai além daqueles previstos nos momentos de parque? - Os espaços favorecem às \textbf{brincadeiras} das \textbf{crianças}, possibilitando : - autonomia da \textbf{criança} na escolha e organização dos parceiros de \textbf{brincar}? - diversidade e contínua inclusão de novos materiais e cantos que ampliem a imaginação? - oportunidades de \textbf{brincar} de forma livre ou mediada pelo professor? - A equipe continuamente se apropria de novas \textbf{brincadeiras} e jogos, enriquecendo o repertório que circula na unidade?
As \textbf{brincadeiras} e jogos tradicionais da comunidade têm sido incluídos no cotidiano de \textbf{bebês} e \textbf{crianças}? - As \textbf{crianças} são \textbf{apoiadas} na resolução de conflitos e problemas que emergem nas situações de \textbf{brincadeira}? - As possibilidades de \textbf{brincadeiras} e jogos com \textbf{bebês}, \textbf{crianças} bem \textbf{pequenas} e \textbf{crianças} \textbf{pequenas} são estudadas e garantidas nos ambientes de aprendizagem da Educação \textbf{Infantil}?
Como se configuram?
Participar ativamente, com \textbf{adultos} e outras \textbf{crianças}, tanto do planejamento da gestão da escola e das atividades propostas pelo educador quanto da realização das atividades da vida cotidiana, tais como a escolha das \textbf{brincadeiras}, dos materiais e dos ambientes, desenvolvendo diferentes linguagens e elaborando conhecimentos, decidindo e se posicionando. - As \textbf{crianças} são convidadas a opinar, por meio de diferentes linguagens ( conversas, desenhos, \textbf{gestos} e movimentos etc. ), sobre a qualidade do tempo que passam na escola?
Sua opinião é considerada? - Como a equipe percebe e favorece o bem-estar e o envolvimento dos \textbf{bebês} ( que ainda não falam ) nos tempos de sua jornada na instituição? - Momentos de deliberação coletiva acerca de definições do planejamento pedagógico são partilhados com as \textbf{crianças}? - As \textbf{crianças} participam da construção da documentação que seguirá com elas para os níveis seguintes de sua escolaridade?
Escolhem as produções mais significativas?
Têm suas elaborações mais pessoais registradas pelo professor?
Explorar movimentos, \textbf{gestos}, \textbf{sons}, formas, \textbf{texturas}, cores, palavras, emoções, transformações, relacionamentos, histórias, objetos, elementos da natureza, na escola e fora dela, ampliando seus saberes sobre a cultura, em suas diversas modalidades : as artes, a escrita, a ciência e a tecnologia. - As \textbf{crianças} têm tido tempo para empreender suas investigações e explorações sobre os mais diferentes objetos e assuntos? - A \textbf{professora}, o professor tem ampliado o seu repertório de investigações e possibilidades de expressão, incluindo novas informações, chamando atenção para fenômenos físicos e sociais e ainda trazendo para o espaço da unidade elementos da contemporaneidade que sejam de interesse ou despertem o interesse da \textbf{criança} por conhecer? - As possibilidades de \textbf{movimentação} corporal e de expressividade gestual da \textbf{criança} têm sido garantidas?
Como?
Em que tempos e espaços?
Expressar, como sujeito dialógico, criativo e sensível, suas necessidades, emoções, sentimentos, dúvidas, hipóteses, descobertas, opiniões, questionamentos, por meio de diferentes linguagens. - \textbf{Bebês} e \textbf{crianças} têm tido espaço para a expressão de suas opiniões, sentimentos e emoções, sendo \textbf{apoiados} na direção da compreensão destes e da interação com o outro? - As hipóteses das \textbf{crianças} são consideradas como pontos de partida para as problematizações e investigações da própria \textbf{criança} e, a depender, do grupo? - As professoras, os professores \textbf{apoiam} as \textbf{crianças} na construção de cenários de suas \textbf{brincadeiras}, alimentando essa construção com novos elementos e garantindo tempo para a produção criativa das \textbf{crianças}? - Literatura \textbf{infantil}, desenhos, \textbf{pinturas}, esculturas, fotografia, arquitetura, dentre outras produções artístico-culturais, fazem parte significativa do cotidiano de \textbf{bebês} e \textbf{crianças} tanto como elementos de fruição e investigação quanto como possibilidades expressivas? - Tecnologias digitais estão a favor da expressão e apreciação de conteúdos que enriqueçam o repertório da \textbf{criança}, superando conteúdo das mídias de massa e do consumismo?
Como?
Em que momentos?
Conhecer -se e construir sua identidade pessoal, social e cultural, constituindo uma imagem positiva de si e de seus grupos de pertencimento, nas diversas experiências de \textbf{cuidados}, interações, \textbf{brincadeiras} e linguagens vivenciadas na instituição escolar e em seu contexto familiar e comunitário. - Situações de preconceito, opressão e discriminação são trabalhadas na instituição, visando sua superação?
Como? \textbf{Bebês} e \textbf{crianças} são protegidos pelos \textbf{adultos} nesse processo? - As referências culturais, históricas, linguísticas, regionais da comunidade escolar são consideradas como referências na construção do projeto pedagógico da instituição?
Como as diversas famílias participam desse processo?
Como \textbf{bebês} e \textbf{crianças} são envolvidos nessa construção? - Professoras, professores constroem critérios de seleção de livros a serem lidos com, pelas e para as \textbf{crianças}, com vistas a ampliar as possibilidades de construção da identidade pelo reconhecimento das narrativas de vida de outras pessoas e personagens?
Essas perguntas — junto com muitas outras que devem emergir no contexto das equipes escolares e dos projetos pedagógicos — nos ajudam a pensar que garantir Direitos de Aprendizagem e Desenvolvimento envolve concretizar esses direitos nos ambientes de aprendizagem que são planejados e replanejados continuamente nas instituições.
É a concretização desses direitos que favorece a construção de uma Educação \textbf{Infantil} democrática e conectada com o desenvolvimento integral de \textbf{bebês} e \textbf{crianças}.
Assim, a oportunidade de ampliação dos estudos sobre Educação \textbf{Infantil} e de participação em experiências pedagógicas de qualidade coopera para elaboração coletiva de Propostas Pedagógicas que respeitem as especificidades das \textbf{crianças} e de cada infância e que podem contribuir para a aprendizagem e desenvolvimento, ampliando as múltiplas linguagens dos \textbf{bebês} e das \textbf{crianças}, tornando -as mais humanas e solidárias.

% matched lemmas: adulto, apoiar, bebê, brincadeira, brincar, criança, cuidado, gesto, infantil, movimentação, pequeno, pintura, professora, sensorial, som, textura
\end{textsample}
